Tres càrregues elèctriques $q_1 = 1,00\ \mathrm{\mu C}$, $q_2 = 3,00\ \mathrm{\mu C}$ i $q_3 = 12,00\ \mathrm{\mu C}$ estan fixades en tres dels vèrtexs del rectangle, tal com es veu en la figura. La distància $d_1$ és de 2,00~m i la distància $d_2$ és de 4,00~m.

\figura[0.45]{fig1.png}

\begin{apartat}{1}
Representeu en un esquema les forces elèctriques que actuen sobre la càrrega $q_1$ per efecte de les altres dues càrregues. Representeu-hi també la força total i calculeu-ne el mòdul.
\end{apartat}

\begin{apartat}{1}
Calculeu el potencial elèctric en el punt $P$ i l'energia potencial de la distribució de les tres càrregues.
\end{apartat}

\blocetiquetat{Dada:}{$k = \dfrac{1}{4\pi\varepsilon_0} = 8,99\times 10^{9}\ \mathrm{N\,m^2\,C^{-2}}$.}
